\section{AIME I 1985}
        \begin{enumerate}
        	\item (\textit{AIME I 1985}) Let $x_1=97$, and for $n>1$ let $x_n=\frac{n}{x_{n-1}}$. Calculate the product $x_1x_2 \ldots x_8$.


 
 \textbf{Answer: }384



\textbf{Solution}

Since $x_n=\frac{n}{x_{n-1}}$, $x_n \cdot x_{n - 1} = n$.  Setting $n = 2, 4, 6$ and $8$ in this equation gives us respectively $x_1x_2 = 2$, $x_3x_4 = 4$, $x_5x_6 = 6$ and $x_7x_8 = 8$ so 
 \[x_1x_2x_3x_4x_5x_6x_7x_8 = 2\cdot4\cdot6\cdot8 = \boxed{384}.\] Notice that the value of $x_1$ was completely unneeded!

 \textbf{Another Way to think about the solution} Every time you multiply in the next term of the sequence, all the numbers before are flipped from the numerator to the denominator or the denominator to the numerator because they are divided. So, 97, the first term, will appear in the multiplied out form in this pattern: $NDNDNDND$ where $N$ is the numerator and $D$ is the denominator. The second term (2) will appear in the pattern $XNDNDNDN$ where $X$ means that the number is skipped the first term. All the pairs of $ND$ cross out and you find that only the even terms have an $N$ left over while all the values of the odd terms are crossed out.

 





\textbf{Solution 2}

Another way to do this is to realize that most of our numbers will be canceled out in the multiplication in the end, and to just list out the terms of our product and cancel:

 
 \[x_1x_2x_3x_4x_5x_6x_7x_8=x_1\cdot\dfrac{2}{x_1}\cdot\dfrac{3}{\dfrac{2}{x_1}}\cdot\dfrac{4}{\dfrac{3}{\dfrac{2}{x_1}}}\cdot\dfrac{5}{\dfrac{4}{\dfrac{3}{\dfrac{2}{x_1}}}}\cdot\dfrac{6}{\dfrac{5}{\dfrac{4}{\dfrac{3}{\dfrac{2}{x_1}}}}}\cdot\dfrac{7}{\dfrac{6}{\dfrac{5}{\dfrac{4}{\dfrac{3}{\dfrac{2}{x_1}}}}}}\cdot\dfrac{8}{\dfrac{7}{\dfrac{6}{\dfrac{5}{\dfrac{4}{\dfrac{3}{\dfrac{2}{x_1}}}}}}}\]
 \[=\left (x_1\cdot\dfrac{2}{x_1} \right )\cdot \left (\dfrac{3}{\dfrac{2}{x_1}}\cdot\dfrac{4}{\dfrac{3}{\dfrac{2}{x_1}}} \right )\cdot \left (\dfrac{5}{\dfrac{4}{\dfrac{3}{\dfrac{2}{x_1}}}}\cdot\dfrac{6}{\dfrac{5}{\dfrac{4}{\dfrac{3}{\dfrac{2}{x_1}}}}} \right )\cdot \left (\dfrac{7}{\dfrac{6}{\dfrac{5}{\dfrac{4}{\dfrac{3}{\dfrac{2}{x_1}}}}}}\cdot\dfrac{8}{\dfrac{7}{\dfrac{6}{\dfrac{5}{\dfrac{4}{\dfrac{3}{\dfrac{2}{x_1}}}}}}} \right )\]
 \[=(2)\cdot (4)\cdot (6)\cdot (8)=\boxed{384}\]

 




	\item (\textit{AIME I 1985}) When a right triangle is rotated about one leg, the volume of the cone produced is $800\pi \;\textrm{cm}^3$. When the triangle is rotated about the other leg, the volume of the cone produced is $1920\pi \;\textrm{cm}^3$. What is the length (in cm) of the hypotenuse of the triangle?


 
 \textbf{Answer: }026



\textbf{Solution}

Let one leg of the triangle have length $a$ and let the other leg have length $b$.  When we rotate around the leg of length $a$, the result is a cone of height $a$ and \href{https://artofproblemsolving.com/wiki/index.php?title=Radius}{radius} $b$, and so of volume $\frac 13 \pi ab^2 = 800\pi$.  Likewise, when we rotate around the leg of length $b$ we get a cone of height $b$ and radius $a$ and so of volume $\frac13 \pi b a^2 = 1920 \pi$.  If we divide this equation by the previous one, we get $\frac ab = \frac{\frac13 \pi b a^2}{\frac 13 \pi ab^2} = \frac{1920}{800} = \frac{12}{5}$, so $a = \frac{12}{5}b$.  Then $\frac{1}{3} \pi \left(\frac{12}{5}b\right)b^2 = 800\pi$ so $b^3 = 1000$ and $b = 10$ so $a = 24$.  Then by the \href{https://artofproblemsolving.com/wiki/index.php?title=Pythagorean\_Theorem}{Pythagorean Theorem}, the hypotenuse has length $\sqrt{a^2 + b^2} = \boxed{026}$.

 


 





\textbf{Solution 2}

Let $a$, $b$ be the $2$ legs, we have the $2$ equations
 \[\frac{a^2b\pi}{3}=800\pi,\frac{ab^2\pi}{3}=1920\pi\]Thus $a^2b=2400, ab^2=5760$. Multiplying gets
 \begin{align*} (a^2b)(ab^2)&=2400\cdot5760\\ (ab)^3&=(2^5\cdot3\cdot5^2)(2^7\cdot3^2\cdot5)\\ ab&=\sqrt[3]{2^{12}\cdot3^3\cdot5^3}=240 \end{align*}
Adding gets
 \begin{align*} a^2b+ab^2=ab(a+b)&=2400+5760\\ 240(a+b)&=240\cdot(10+24)\\ a+b&=34 \end{align*}
Let $h$ be the hypotenuse then
 \begin{align*} h&=\sqrt{a^2+b^2}\\ &=\sqrt{(a+b)^2-2ab}\\ &=\sqrt{34^2-2\cdot240}\\ &=\sqrt{676}\\ &=\boxed{26} \end{align*}


 





\textbf{Solution 3(Ratios)}

Let $a$ and $b$ be the two legs of the equation. We can find $\frac{a}{b}$ by doing $\frac{1920\pi}{800\pi}$. This simplified is $\frac{12}{5}$. We can represent the two legs as $12x$ and $5x$ for $a$ and $b$ respectively. 

 Since the volume of the first cone is $800\pi$, we use the formula for the volume of a cone and get $100\pi x^3=800 \pi$. Solving for $x$, we get $x=2$. 

 Plugging in the side lengths to the Pythagorean Theorem, we get an answer of $\boxed{026}$.

 





\textbf{Solution 4 (Smarter Ratios)}

Continue solution 3 until we obtain, $\frac{x}{y} = \frac{12}{5}$. Let $x = 12$ and $y=5$, and $xy^2 = 300$. This is 8 times off. So, we wish to give a 2 to $x$ and a 4 to $y^2$, meaning that both $x$ and $y$ obtain a factor 2. Thus, $x= 12\cdot 2 = 24$ and $y = 5 \cdot 2 = 10$, and the result is just $\sqrt{676} = \boxed{026}$

 




	\item (\textit{AIME I 1985}) Find $c$ if $a$, $b$, and $c$ are positive integers which satisfy $c=(a + bi)^3 - 107i$, where $i^2 = -1$.


 
 \textbf{Answer: }198



\textbf{Solution}

Expanding out both sides of the given \href{https://artofproblemsolving.com/wiki/index.php?title=Equation}{equation} we have $c + 107i = (a^3 - 3ab^2) + (3a^2b - b^3)i$.  Two \href{https://artofproblemsolving.com/wiki/index.php?title=Complex\_number}{complex numbers} are equal if and only if their \href{https://artofproblemsolving.com/wiki/index.php?title=Real\_part}{real parts} and \href{https://artofproblemsolving.com/wiki/index.php?title=Imaginary\_part}{imaginary parts} are equal, so $c = a^3 - 3ab^2$ and $107 = 3a^2b - b^3 = (3a^2 - b^2)b$.  Since $a, b$ are \href{https://artofproblemsolving.com/wiki/index.php?title=Integer}{integers}, this means $b$ is a \href{https://artofproblemsolving.com/wiki/index.php?title=Divisor}{divisor} of 107, which is a \href{https://artofproblemsolving.com/wiki/index.php?title=Prime\_number}{prime number}.  Thus either $b = 1$ or $b = 107$.  If $b = 107$, $3a^2 - 107^2 = 1$ so $3a^2 = 107^2 + 1$, but $107^2 + 1$ is not divisible by 3, a contradiction.  Thus we must have $b = 1$, $3a^2 = 108$ so $a^2 = 36$ and $a = 6$ (since we know $a$ is positive).  Thus $c = 6^3 - 3\cdot 6 = \boxed{198}$.

 \textbf{Video Solution by SpreadTheMathLove} 




	\item (\textit{AIME I 1985}) A small  square is constructed inside a square of area $1$ by dividing each side of the unit square into $n$ equal parts, and then connecting the  vertices to the division points closest to the opposite vertices, as shown in the figure. Find the value of $n$ if the area of the small square is exactly $\frac1{1985}$.


 
                    \begin{figure}[H]
                        \centering
                        \includegraphics[width=0.5\textwidth]{Problem/AIME_1985_I_4_Problem_Figure_1.png}
                    \end{figure}
                    


 
 \textbf{Answer: }032



\textbf{Solution 1}

The lines passing through $A$ and $C$ divide the square into three parts, two \href{https://artofproblemsolving.com/wiki/index.php?title=Right\_triangle}{right triangles} and a \href{https://artofproblemsolving.com/wiki/index.php?title=Parallelogram}{parallelogram}.  Using the smaller side of the parallelogram, $1/n$, as the base, where the height is 1, we find that the area of the parallelogram is $A = \frac{1}{n}$.  By the \href{https://artofproblemsolving.com/wiki/index.php?title=Pythagorean\_Theorem}{Pythagorean Theorem}, the longer base of the parallelogram has \href{https://artofproblemsolving.com/wiki/index.php?title=Length}{length} $l = \sqrt{1^2 + \left(\frac{n - 1}{n}\right)^2} = \frac{1}{n}\sqrt{2n^2 - 2n + 1}$, so the parallelogram has height $h = \frac{A}{l} = \frac{1}{\sqrt{2n^2 - 2n + 1}}$.  But the height of the parallelogram is the side of the little square, so $2n^2 - 2n + 1 = 1985$.  Solving this \href{https://artofproblemsolving.com/wiki/index.php?title=Quadratic\_equation}{quadratic equation} gives $n = \boxed{32}$.

 





\textbf{Solution 2}


                    \begin{figure}[H]
                        \centering
                        \includegraphics[width=0.5\textwidth]{Solution/AIME_1985_I_4_Solution_Figure_1.png}
                    \end{figure}
                    

 Surrounding the square with area $\frac{1}{1985}$ are $4$ right triangles with hypotenuse $1$ (sides of the large square). Thus, $X + \frac{1}{1985} = 1$, where $X$ is the area of the of the 4 triangles.We can thus use proportions to solve this problem.
 \begin{eqnarray*} \frac{GF}{BE}=\frac{CG}{CB}\implies \frac{\frac{1}{\sqrt{1985}}}{BE}=\frac{\frac{1}{n}}{1}\implies BE=\frac{n\sqrt{1985}}{1985} \end{eqnarray*}Also,
 \begin{eqnarray*} \frac{BE}{1}=\frac{EC}{\frac{n-1}{n}}\implies EC=\frac{\sqrt{1985}}{1985}(n-1) \end{eqnarray*}Thus,
 \begin{eqnarray*} 2(BE)(EC)+\frac{1}{1985}=1\\ 2n^{2}-2n+1=1985\\ n(n-1)=992 \end{eqnarray*}Simple factorization and guess and check gives us $\boxed{32}$.

 





\textbf{Solution 3}


                    \begin{figure}[H]
                        \centering
                        \includegraphics[width=0.5\textwidth]{Solution/AIME_1985_I_4_Solution_Figure_1.png}
                    \end{figure}
                    

 Line Segment $DE = \frac{1}{n}$, so $EC = 1 - \frac{1}{n} = \frac{n-1}{n}$. Draw line segment $HE$ parallel to the corresponding sides of the small square, $HE$ has length $\frac{1}{\sqrt{1985}}$, as it is the same length as the sides of the square. Notice that $\triangle CEL$ is similar to $\triangle HDE$ by $AA$ similarity. Thus, $\frac{LC}{HE} = \frac{EC}{DE} = n-1$, so $LC = \frac{n-1}{\sqrt{1985}}$. Notice that $\triangle CEL$ is also similar to $\triangle CDF$ by $AA$ similarity. Thus, $\frac{FC}{EC} = \frac{DC}{LC}$, and the expression simplifies into a quadratic equation $n^2 - n - 992 = 0$. Solving this quadratic equation yields $n =\boxed{32}$.

 




	\item (\textit{AIME I 1985}) A sequence of integers $a_1, a_2, a_3, \ldots$ is chosen so that $a_n = a_{n - 1} - a_{n - 2}$ for each $n \ge 3$. What is the sum of the first $2001$ terms of this sequence if the sum of the first $1492$ terms is $1985$, and the sum of the first $1985$ terms is $1492$?


 
 \textbf{Answer: }986



\textbf{Solution 1}

The problem gives us a sequence defined by a \href{https://artofproblemsolving.com/wiki/index.php?title=Recursion}{recursion}, so let's calculate a few values to get a feel for how it acts.  We aren't given initial values, so let $a_1 = a$ and $a_2 = b$.  Then $a_3 = b - a$, $a_4 = (b - a) - b = -a$, $a_5 = -a - (b - a) = -b$, $a_6 = -b - (-a) = a - b$, $a_7 = (a - b) - (-b) = a$ and $a_8 = a - (a - b) = b$.  Since the sequence is recursively defined by the first 2 terms, after this point it must continue to repeat.  Thus, in particular $a_{j + 6} = a_j$ for all $j$, and so repeating this $n$ times, $a_{j + 6n} = a_j$ for all \href{https://artofproblemsolving.com/wiki/index.php?title=Integer}{integers} $n$ and $j$.  

 Because of this, the sum of the first 1492 terms can be greatly simplified: $1488 = 6 \cdot 248$ is the largest multiple of 6 less than 1492, so 
 \[\sum_{i = 1}^{1492} a_i = (a_{1489} + a_{1490} + a_{1491} + a_{1492})+ \sum_{i = 1}^{1488} a_i = (a_1 + a_2 + a_3 + a_4) + \sum_{n = 0}^{247}\sum_{j = 1}^6 a_{6n + j}\]
 \[=(a + b + (b - a) + (-a)) + \sum_{n = 0}^{247}\sum_{j = 1}^6 a_j = 2b - a,\]where we can make this last step because $\sum_{j = 1}^6 a_j = 0$ and so the entire second term of our \href{https://artofproblemsolving.com/wiki/index.php?title=Expression}{expression} is \href{https://artofproblemsolving.com/wiki/index.php?title=0}{zero}.

 Similarly, since $1980 = 6 \cdot 330$, $\sum_{i = 1}^{1985} a_i = (a_1 + a_2 + a_3 + a_4 + a_5) + \sum_{i = 1}^{1980}a_i = a + b + (b - a) + (-a) + (-b) = b - a$.

 Finally, $\sum_{i = 1}^{2001}a_i = a_1 + a_2 + a_3 + \sum_{i = 1}^{1998} a_i = a + b + (b - a) = 2b$.

 Then by the givens, $2b - a = 1985$ and $b - a = 1492$ so $b = 1985 - 1492 = 493$ and so the answer is $2\cdot 493 = \boxed{986}$.

 





\textbf{Solution 2 (Beginner Friendly)}

We will start by letting the first 2 numbers in the sequence be $A$ and $B$. 

 Then the sequence would go $A$, $B$, $B$ - $A$, -$A$, -$B$, $A$ - $B$, $A$, $B$ $....$

 Here we can see that the sequence repeats itself every 6 terms.

 If we add up all of those 6 terms we will get 0.

 Now we are given the sum of the first 1985 terms is 1492 and the sum of the first 1492 terms is 1985. 

 If the first 20 terms = $X$ then we also can say that $term 19$ + $term 20$ = $X$. Since every group of 6 will be equal to 0.

 The first 1985 terms sum to 1492 $\longrightarrow$ terms 1981, 1982, 1983, 1984, and 1985 sum to 1492.

 So $A$ + $B$ + $B$ - $A$ - $A$ - $B$ = 1492 $\longrightarrow$ $B$ - $A$ = 1492

 The first 1492 terms sum to 1985 $\longrightarrow$ terms 1489, 1490, 1491, and 1492 sum to 1985

 So $A$ + $B$ + $B$ - $A$ - $A$ = 1985 $\longrightarrow$ $2B$ - $A$ = 1985

 Now we solve given:  [$2B$ - $A$ = 1985] and [$B$ - $A$ = 1492]

 [$2B$ - $A$] - [$B$ - $A$] = 1985 - 1492 $\longrightarrow$ $B$ = 493

 [$B$ - $A$] = 1492 $\longrightarrow$ [$493$ - $A$] = 1492 $\longrightarrow$ $A$ = -999

 Now we need to find the sum of the first 2001 terms, which is the same as the sum of terms 1999, 2000, and 2001.

 So our answer is $A$ + $B$ + $B$ - $A$ $\longrightarrow$ $2B$ $\longrightarrow$ $2\cdot493$ = $\boxed{986}$.

 




	\item (\textit{AIME I 1985}) As shown in the figure, $\triangle ABC$ is divided into six smaller triangles by lines drawn from the vertices through a common interior point. The areas of four of these triangles are as indicated. Find the area of $\triangle ABC$.


 
                    \begin{figure}[H]
                        \centering
                        \includegraphics[width=0.5\textwidth]{Problem/AIME_1985_I_6_Problem_Figure_1.png}
                    \end{figure}
                    


 
 \textbf{Answer: }315



\textbf{Solution 1}

Let the interior point be $P$, let the points on $\overline{BC}$, $\overline{CA}$ and $\overline{AB}$ be $D$, $E$ and $F$, respectively.  Let $x$ be the area of $\triangle APE$ and $y$ be the area of $\triangle CPD$.  Note that $\triangle APF$ and $\triangle BPF$ share the same \href{https://artofproblemsolving.com/wiki/index.php?title=Altitude}{altitude} from $P$, so the \href{https://artofproblemsolving.com/wiki/index.php?title=Ratio}{ratio} of their areas is the same as the ratio of their bases.  Similarly, $\triangle ACF$ and $\triangle BCF$ share the same altitude from $C$, so the ratio of their areas is the same as the ratio of their bases.  Moreover, the two pairs of bases are actually the same, and thus in the same ratio.  As a result, we have:$\frac{40}{30} = \frac{124 + x}{65 + y}$ or equivalently $372 + 3x = 260 + 4y$ and so $4y = 3x+ 112$.

 Applying identical reasoning to the triangles with bases $\overline{CD}$ and $\overline{BD}$, we get $\frac{y}{35} = \frac{x+y+84}{105}$ so that $3y = x + y + 84$ and $2y = x + 84$.  Substituting from this equation into the previous one gives $x = 56$, from which we get $y = 70$ and so the area of $\triangle ABC$ is $56 + 40 + 30 + 35 + 70 + 84 = \Rightarrow \boxed{315}$.

 





\textbf{Solution 2 (Mass Points)}


                    \begin{figure}[H]
                        \centering
                        \includegraphics[width=0.5\textwidth]{Solution/AIME_1985_I_6_Solution_Figure_1.png}
                    \end{figure}
                    We can solve this problem using \href{https://artofproblemsolving.com/wiki/index.php?title=Mass\_points}{mass points}. First, denote the interior point as $P$, the foot of the \href{https://artofproblemsolving.com/wiki/index.php?title=Cevian}{cevian} from $A$ as $A'$, the foot of the cevian from $B$ as $B'$, and the foot of the cevian from $C$ as $C'$. Notice how triangles $\Delta APC'$ and $\Delta BPC'$ share the same \href{https://artofproblemsolving.com/wiki/index.php?title=Altitude}{altitude} from $P$. Because their areas are in the ratio $40:30=4:3$, then their bases $AC'$ and $BC'$ must also be in the ratio $4:3$. Next, notice how triangles $\Delta APB$ and $\Delta PBA'$ also share the same altitude from B due to the fact A,P,A' are collinear. Their areas are in the ratio $40+30:35=70:35=2:1$, so the ratio $AP:PA'$ must also be $2:1$.

 Now, we can apply mass points. First, assign vertex $B$ a weight of $1$. Then, we determine the weight of $A$ as $\frac{BC'}{AC'}\cdot W_{B}=\frac{3}{4}$, and by adding the weights of $A$ and $B$, we get $W_{C'}=\frac{7}{4}$. Applying the same to $\overline{AA'}$, we get that $W_{A'}=\frac{3}{2}$, $W_{P}=\frac{9}{4}$, $W_{C}=\frac{1}{2}$, and $W_{B'}=\frac{5}{4}$.

 Finally, we can determine the area of $\Delta ABC$ by adding together the areas of all of the small triangles, and using side length ratios derived from mass points to determine the areas of the two unknown triangles.
 \[[\Delta ABC]=40+30+35+84+[\Delta APB']+[\Delta CPA']\]
 \[[\Delta ABC]=189+\frac{2}{3}\cdot [\Delta CPB']+2\cdot [\Delta BPA']\]
 \[[\Delta ABC]=189+56+70\]
 \[[\Delta ABC]=\boxed{315}\]

 





\textbf{Solution 3}

Let the interior point be $P$ and let the points on $\overline{BC}$, $\overline{CA}$ and $\overline{AB}$ be $D$, $E$ and $F$, respectively. Also, let $[APE]=x,[CPD]=y.$ Then notice that by Ceva's, $\frac{FB\cdot DC\cdot EA}{DB\cdot CE\cdot AF}=1.$ However, we can deduce $\frac{FB}{AF}=\frac{3}{4}$ from the fact that $[AFP]$ and $[BPF]$ share the same height. Similarly, $x=\frac{84CE}{EA}$ and $y=\frac{35DC}{BD}.$ Plug and chug and you get $xy=84\cdot 35\cdot \frac{3}{4}=2205.$ Then notice by the same height reasoning, $\frac{84}{x}=\frac{119+y}{x+70}.$ Clear the fractions and combine like terms to get $35x=5880-xy.$ We know $xy=2205$ so subtraction yields $35x=3675,$ or $x=105.$ Plugging this in to our previous ratio statement yields $\frac{84}{105}=\frac{4}{5}=\frac{119+y}{175},$ so $y=21.$ Basic addition gives us $105+84+21+35+30+40=\boxed{315}.$

 -dchen

 





\textbf{Solution 4}

Let the interior point be $P$ and let the points on $\overline{BC}$, $\overline{CA}$ and $\overline{AB}$ be $D$, $E$ and $F$, respectively. Then the cevians $AD,BF,CE$ are concurrent, so we can use Ceva's Theorem, letting $\frac{BD}{DC}=\frac{a}{b}$ and $\frac{CF}{FA}=\frac{c}{d}$. Notice that $\frac{AE}{EB}=\frac{[\Delta APE]}{[\Delta EPB]}=\frac43.$ 
 \[\frac{4}{3}\cdot \frac{a}{b}\cdot \frac{c}{d}=1\implies \frac{d}{c}=\frac{a}{b}\cdot \frac{4}{3}.\]

 We know that $[\Delta CPD]=35\cdot \frac ba$ and $[\Delta APF]=84\cdot \frac dc,$ so
 \[[\Delta ABC] = 84+84\cdot \frac dc + 35\cdot \frac ba + 40+30+35 = \left(1+\frac ba\right)(40+30+35).\]We will now solve for $\frac ba$:

 
 \[84+84\cdot \frac ab\cdot \frac 43 + 35\cdot \frac ba = \frac ba\cdot 105.\]
 \[12+16\cdot \frac ab + 5\cdot \frac ba = \frac ba\cdot 15\]
 \[10\left(\frac ba\right)^2-12\cdot \frac ba-16=0\]Factoring this gives $\left(\frac ba-2\right)\left(10\cdot \frac ba - 8\right)=0,$ so the area of $\triangle ABC$ is
 \[\left(1+\frac ba\right)(40+30+35)=3\cdot 105=\boxed{315.}\]

 





\textbf{Solution 5 (Ceva Theorem, Height to Base Relationships)}

That common interior point screams Ceva’s theorem. However, we have no idea how to apply it yet! Let the common interior point be $P$. It follows that if the cevians from A, B, C, to the opposite sides intersect at $D,E,F$ respectively, then triangle $APF$ and $FPB$ share a height, meaning the ratio of their bases is the ratio of their areas, or $\frac{40}{30}$. Let the area of triangles $APE$ and $PCD$ be $x$ and $y$. Then, notice that triangles $AFC$ and $FCB$ share a height, so the ratio of their areas is the same as the ratio of their bases. The ratio of their bases however is $\frac{40}{30}$, so the ratio of their areas is $\frac{124+x}{65+y} = \frac{40}{30}$. Now understand that the ratio of the remaining sides is $\frac{35}{y}$ and $\frac{84}{x}$, so through Ceva’s Theorem, $\frac{4 \cdot 35 \cdot 28}{ x\cdot y} = 1$, so $xy = 3920$. From the above relation we know that $\frac{124+x}{65+y} = \frac{40}{30}$, so a tad bit of algebra gives us two possible ordered pairs $(x,y)$, either $(56,70)$, or $(-\frac{288}{3}, -42)$. Obviously, the negative area is nonexistent in Plane Euclidean Geometry. Therefore, the first ordered pair is true, and the area is just $56+70+84+40+30+35 = \boxed{315}$.

 \textbf{Note} This problem gives redundant information. The area of any one of the 4 triangles given can be derived from the other three triangles' areas

 




	\item (\textit{AIME I 1985}) Assume that $a$, $b$, $c$ and $d$ are positive integers such that $a^5 = b^4$, $c^3 = d^2$ and $c - a = 19$. Determine $d - b$.


 
 \textbf{Answer: }757



\textbf{Solution}

It follows from the givens that $a$ is a \href{https://artofproblemsolving.com/wiki/index.php?title=Perfect\_power}{ perfect fourth power}, $b$ is a perfect fifth power, $c$ is a \href{https://artofproblemsolving.com/wiki/index.php?title=Perfect\_square}{perfect square} and $d$ is a \href{https://artofproblemsolving.com/wiki/index.php?title=Perfect\_cube}{perfect cube}.  Thus, there exist \href{https://artofproblemsolving.com/wiki/index.php?title=Integer}{integers} $s$ and $t$ such that $a = t^4$, $b = t^5$, $c = s^2$ and $d = s^3$.  So $s^2 - t^4 = 19$.  We can factor the left-hand side of this \href{https://artofproblemsolving.com/wiki/index.php?title=Equation}{equation} as a difference of two squares, $(s - t^2)(s + t^2) = 19$.  $19$ is a \href{https://artofproblemsolving.com/wiki/index.php?title=Prime\_number}{prime number} and $s + t^2 > s - t^2$ so we must have $s + t^2 = 19$ and $s - t^2 = 1$.  Then $s = 10, t = 3$ and so $d = s^3 = 1000$, $b = t^5 = 243$ and $d-b=\boxed{757}$.

 




	\item (\textit{AIME I 1985}) The sum of the following seven numbers is exactly 19: $a_1 = 2.56$, $a_2 = 2.61$, $a_3 = 2.65$, $a_4 = 2.71$, $a_5 = 2.79$, $a_6 = 2.82$, $a_7 = 2.86$. It is desired to replace each $a_i$ by an integer approximation $A_i$, $1\le i \le 7$, so that the sum of the $A_i$'s is also $19$ and so that $M$, the maximum of the "errors" $\lvert A_i-a_i \rvert$, is as small as possible. For this minimum $M$, what is $100M$?


 
 \textbf{Answer: }061



\textbf{Solution 1}

Let \( \{x\} \) denote the fractional part of \( x \).

 Because \( \{2.86\} - \{2.56\} = 0.3 \), we see that the maximum spread of the numbers is \( 0.3 \), and the median spread of the numbers is \( 0.15 \).

 Now, we take inspiration from rounding decimals. Note how if we have one positive integer \( a \) and another positive integer \( a+1 \), then \( a+1-a=1 \), and the median spread is \( 1/2 \), or \( 0.5 \). Then, the following piecewise function is acquired. If a number \( n \) satisfies \( a < n < a+1 \), then if \( \{n\} < 0.5 \), we take \( \lfloor{n} \rfloor \), and if \( \{n\} \ge 0.5 \), we take \( \lceil{n} \rceil \).

 Applying this to our scenario, we have \( 2 < a_i < 3 \) and the median spread is \( 2.5 \). So, to determine a judgement in our scenario, we use the following piecewise function of \( A_i \).

 If \( \{a_i\} < 2.5 + 0.15 = 2.65 \), then take \( \lfloor{a_i} \rfloor \), and if \( \{a_i\} \ge 2.65 \), take \( \lceil{a_i} \rceil \).

 For all 7 numbers, we apply these and get 2, 2, 3, 3, 3, 3, 3. Adding them results in \( 3 \cdot 5 + 4 = 19 \), which satisfies the constraint, and the greatest value of \( M \) is determined from \( 2.61 - \lfloor{2.61} \rfloor = 0.61 \).

 This is guarenteed as notice that all \( a_i > 2.5 \), and in general the rounding principle states that if a number \( n \) has the property that \( \{n\} < \{2.5\} \), then take \( \lfloor{n} \rfloor \) and if \( \{n\} \ge \{2.5\} \), then take \( \lceil{n} \rceil \), in which we can adjust using the same principle here.

 In addition to this problem, this idea is reminiscent for ALL problems of this category, so one additional concept to add to your problem solving toolkit.

 Thus, \( M = 0.61 \), and \( 100M = 100 \times 0.61 = \)$\boxed{061}$

 





\textbf{Solution 2}

If any of the approximations $A_i$ is less than 2 or more than 3, the error associated with that term will be larger than 1, so the largest error will be larger than 1.  However, if all of the $A_i$ are 2 or 3, the largest error will be less than 1.  So in the best case, we write 19 as a sum of 7 numbers, each of which is 2 or 3.  Then there must be five 3s and two 2s.  It is clear that in the best appoximation, the two 2s will be used to approximate the two smallest of the $a_i$, so our approximations are $A_1 = A_2 = 2$ and $A_3 = A_4 = A_5 = A_6 = A_7 = 3$ and the largest error is $|A_2 - a_2| = 0.61$, so the answer is $\boxed{061}$.

 




	\item (\textit{AIME I 1985}) In a circle, parallel chords of lengths $2$, $3$, and $4$ determine central angles of $\alpha$, $\beta$, and $\alpha + \beta$ radians, respectively, where $\alpha + \beta < \pi$. If $\cos \alpha$, which is a positive rational number, is expressed as a fraction in lowest terms, what is the sum of its numerator and denominator?


 
 \textbf{Answer: }049



\textbf{Solution 1}


                    \begin{figure}[H]
                        \centering
                        \includegraphics[width=0.5\textwidth]{Solution/AIME_1985_I_9_Solution_Figure_1.png}
                    \end{figure}
                    All chords of a given length in a given circle subtend the same \href{https://artofproblemsolving.com/wiki/index.php?title=Arc}{arc} and therefore the same central angle.  Thus, by the given, we can re-arrange our chords into a \href{https://artofproblemsolving.com/wiki/index.php?title=Triangle}{triangle} with the circle as its \href{https://artofproblemsolving.com/wiki/index.php?title=Circumcircle}{circumcircle}.  

 
                    \begin{figure}[H]
                        \centering
                        \includegraphics[width=0.5\textwidth]{Solution/AIME_1985_I_9_Solution_Figure_1.png}
                    \end{figure}
                    This triangle has \href{https://artofproblemsolving.com/wiki/index.php?title=Semiperimeter}{semiperimeter} $\frac{2 + 3 + 4}{2}$ so by \href{https://artofproblemsolving.com/wiki/index.php?title=Heron\%27s\_formula}{Heron's formula} it has \href{https://artofproblemsolving.com/wiki/index.php?title=Area}{area} $K = \sqrt{\frac92 \cdot \frac52 \cdot \frac32 \cdot \frac12} = \frac{3}{4}\sqrt{15}$.  The area of a given triangle with sides of length $a, b, c$ and circumradius of length $R$ is also given by the formula $K = \frac{abc}{4R}$, so $\frac6R = \frac{3}{4}\sqrt{15}$ and $R = \frac8{\sqrt{15}}$.

 Now, consider the triangle formed by two radii and the chord of length 2.  This \href{https://artofproblemsolving.com/wiki/index.php?title=Isosceles\_triangle}{isosceles triangle} has vertex angle $\alpha$, so by the \href{https://artofproblemsolving.com/wiki/index.php?title=Law\_of\_Cosines}{Law of Cosines},

 
 \[2^2 = R^2 + R^2 - 2R^2\cos \alpha \Longrightarrow \cos \alpha = \frac{2R^2 - 4}{2R^2} = \frac{17}{32}\]and the answer is $17 + 32 = \boxed{049}$.

 





\textbf{Solution 2 (Law of Cosines)}


                    \begin{figure}[H]
                        \centering
                        \includegraphics[width=0.5\textwidth]{Solution/AIME_1985_I_9_Solution_Figure_1.png}
                    \end{figure}
                    It’s easy to see in triangle which lengths 2, 3, and 4, that the angle opposite the side 2 is $\frac{\alpha}{2}$, and using the \href{https://artofproblemsolving.com/wiki/index.php?title=Law\_of\_Cosines}{Law of Cosines}, we get:  
 \[2^2 = 3^2 + 4^2 - 2\cdot3\cdot4\cos\frac{\alpha}{2}\] Which, rearranges to: 
 \[21 = 24\cos\frac{\alpha}{2}\] And, that gets us: 
 \[\cos\frac{\alpha}{2} = 7/8\] Using $\cos 2\theta = 2\cos^2 \theta - 1$, we get that: 
 \[\cos\alpha = 17/32\]  Which gives an answer of $\boxed{049}$

 
- AlexLikeMath

 





\textbf{Solution 3 (trig)}

Using the first diagram above,
 \[\sin \frac{\alpha}{2} = \frac{1}{r}\]
 \[\sin \frac{\beta}{2} = \frac{1.5}{r}\]
 \[\sin(\frac{\alpha}{2}+\frac{\beta}{2})=\frac{2}{r}\]by the Pythagorean trig identities,
 \[\cos\frac{\alpha}{2}=\sqrt{1-\frac{1}{r^2}}\]
 \[\cos\frac{\beta}{2}=\sqrt{1-\frac{2.25}{r^2}}\]so by the composite sine identity
 \[\frac{2}{r}=\frac{1}{r}\sqrt{1-\frac{2.25}{r^2}}+\frac{1.5}{r}\sqrt{1-\frac{1}{r^2}}\]multiply both sides by $2r$, then subtract $\sqrt{4-\frac{9}{r^2}}$ from both sidessquaring both sides, we get
 \[16 - 8\sqrt{4-\frac{9}{r^2}} + 4 - \frac{9}{r^2}=9 - \frac{9}{r^2}\]
 \[\Longrightarrow 16+4=9+8\sqrt{4-\frac{9}{r^2}}\Longrightarrow\frac{11}{8}=\sqrt{4-\frac{9}{r^2}}\Longrightarrow\frac{121}{64}=4-\frac{9}{r^2}\]
 \[\Longrightarrow\frac{(256-121)r^2}{64}=9\Longrightarrow r^2= \frac{64}{15}\]plugging this back in,
 \[\cos^2(\frac{\alpha}{2})=1-\frac{15}{64}=\frac{49}{64}\]so
 \[\cos(\alpha)=2(\frac{49}{64})-1=\frac{34}{64}=\frac{17}{32}\]and the answer is $17+32=\boxed{049}$

 




	\item (\textit{AIME I 1985}) How many of the first $1000$ positive integers can be expressed in the form


 $\lfloor 2x \rfloor + \lfloor 4x \rfloor + \lfloor 6x \rfloor + \lfloor 8x \rfloor$,


 where $x$ is a real number, and $\lfloor z \rfloor$ denotes the greatest integer less than or equal to $z$? 


 
 \textbf{Answer: }600



\textbf{Solution 1}

Noting that all of the numbers are even, we can reduce this to any real number $x$ between $0$ to $\frac 12$, as this will be equivalent to $\frac n2$ to $\frac {n+1}2$ for any integer $n$ (same reasoning as above). So now we only need to test every 10 numbers; and our answer will be 100 times the number of integers we can reach between 1 and 10. 

 We can now approach this by directly searching for the integers (this solution) or brute forcing all of the cases (next solution):

 We can match up the greatest integer functions with one of the partitions of the integer. If we let $x = \frac 12$ then we get the solution $10$; now consider when $x < \frac 12$: $\lfloor 2x \rfloor = 0$, $\lfloor 4x \rfloor \le 1$, $\lfloor 6x \rfloor \le 2$, $\lfloor 8x \rfloor \le 3$. But according to this the maximum we can get is $1+2+3 = 6$, so we only need to try the first 6 numbers.

 
                \begin{itemize}
                    \item $1$: Easily possible, for example try plugging in $x =\frac 18$.\item $2$: Also simple, for example using $\frac 16$.\item $3$: The partition must either be $1+1+1$ or $1+2$. If $\lfloor 4x \rfloor = 1$, then $x \ge \frac 14$, but then $\lfloor 8x \rfloor \ge 2$; not possible; and vice versa to show that the latter partition doesn't work. So we cannot obtain $3$.\item $4$: We can partition as $1+1+2$, and from the previous case we see that $\frac 14$ works.\item $5$: We can partition as $1+2+2$, from which we find that $\frac 13$ works.\item $6$: We can partition as $1+2+3$, from which we find that $\frac 38$ works.
                \end{itemize}


                Out of these 6 cases, only 3 fails. So between 1 and 10 we can reach only the integers $1,2,4,5,6,10$; hence our solution is $6 \cdot 100 = \boxed{600}$.

 





\textbf{Solution 2}

As we change the value of $x$, the value of our \href{https://artofproblemsolving.com/wiki/index.php?title=Expression}{expression} changes only when $x$ crosses \href{https://artofproblemsolving.com/wiki/index.php?title=Rational\_number}{rational number} of the form $\frac{m}{n}$, where $n$ is divisible by 2, 4, 6 or 8.  Thus, we need only see what happens at the numbers of the form $\frac{m}{\textrm{lcm}(2, 4, 6, 8)} = \frac{m}{24}$.  This gives us 24 calculations to make; we summarize the results here:

 $\frac{1}{24}, \frac{2}{24} \to 0$

 $\frac{3}{24} \to 1$

 $\frac{4}{24}, \frac{5}{24} \to 2$

 $\frac{6}{24}, \frac{7}{24} \to 4$

 $\frac{8}{24} \to 5$

 $\frac{9}{24}, \frac{10}{24}, \frac{11}{24} \to 6$

 $\frac{12}{24}, \frac{13}{24}, \frac{14}{24} \to 10$

 $\frac{15}{24} \to 11$

 $\frac{16}{24},\frac{17}{24} \to 12$

 $\frac{18}{24}, \frac{19}{24} \to 14$

 $\frac{20}{24}\to 15$

 $\frac{21}{24}, \frac{22}{24}, \frac{23}{24} \to16$

 $\frac{24}{24} \to 20$

 Thus, we hit 12 of the first 20 integers and so we hit $50 \cdot 12 = \boxed{600}$ of the first $1000$.

 





\textbf{Solution 2 Shortcut}

Because $2,4,6,8$ are all multiples of $2$, we can speed things up. We only need to check up to $\frac{12}{24}$, and the rest should repeat. As shown before, we hit 6 integers ($1,2,4,5,6,10$) from $\frac{1}{24}$ to $\frac{12}{24}$. Similarly, this should repeat 100 times, for $\boxed{600}$

 





\textbf{Solution 2 Bigger Shortcut (Close to Solution 6)}

We only need to check the numbers where it increments, namely $\frac{1}{8}, \frac{1}{6}, \frac{1}{4}, \frac{1}{3}, \frac{3}{8}, \frac{1}{2}$. As shown before, we hit 6 integers ($1,2,4,5,6,10$) from $\frac{1}{24}$ to $\frac{1}{2}$. Similarly, this should repeat 100 times, for $\boxed{600}$

 





\textbf{Solution 3}

Recall from Hermite's Identity that $\sum_{k = 0}^{n - 1}\left\lfloor x + \frac kn\right\rfloor = \lfloor nx\rfloor$. Then we can rewrite $\lfloor 2x \rfloor + \lfloor 4x \rfloor + \lfloor 6x \rfloor + \lfloor 8x \rfloor = 4\lfloor x\rfloor + \left\lfloor x + \frac18\right\rfloor + \left\lfloor x + \frac16\right\rfloor + 2\left\lfloor x + \frac14\right\rfloor + \left\lfloor x + \frac13\right\rfloor$$+ \left\lfloor x + \frac38\right\rfloor + 4\left\lfloor x + \frac12\right\rfloor + \left\lfloor x + \frac58\right\rfloor + \left\lfloor x + \frac23\right\rfloor + 2\left\lfloor x + \frac34\right\rfloor + \left\lfloor x + \frac56\right\rfloor + \left\lfloor x + \frac78\right\rfloor$. There are $12$ terms here (we don't actually have to write all of it out; we can just see where there will be duplicates and subtract accordingly from $20$). Starting from every integer $x$, we can keep adding to achieve one higher value for each of these terms, but after raising the last term, we will have raised the whole sum by $20$ while only achieving $12$ of those $20$ values. We can conveniently shift the $1000$ (since it can be achieved) to the position of the $0$ so that there are only complete cycles of $20$, and the answer is $\frac {12}{20}\cdot1000 = \boxed{600}$.

 





\textbf{Solution 4}

Let $x=\lfloor x\rfloor+\{x\}$ then
 \begin{align*} \lfloor 2x\rfloor+\lfloor 4x\rfloor+\lfloor 6x\rfloor+\lfloor 8x\rfloor&=\lfloor 2(\lfloor x\rfloor+\{x\})\rfloor+\lfloor 4(\lfloor x\rfloor+\{x\})\rfloor+\lfloor 6(\lfloor x\rfloor+\{x\})\rfloor+\lfloor 8(\lfloor x\rfloor+\{x\})\rfloor\\ &=2\lfloor x\rfloor+4\lfloor x\rfloor+6\lfloor x\rfloor+8\lfloor x\rfloor+\lfloor 2\{x\}\rfloor+\lfloor 4\{x\}\rfloor+\lfloor 6\{x\}\rfloor+\lfloor 8\{x\}\rfloor\\ &=20\lfloor x\rfloor+(\lfloor 2\{x\}\rfloor+\lfloor 4\{x\}\rfloor+\lfloor 6\{x\}\rfloor+\lfloor 8\{x\}\rfloor) \end{align*}
Similar to the previous solutions, the value of $\lfloor 2\{x\}\rfloor+\lfloor 4\{x\}\rfloor+\lfloor 6\{x\}\rfloor+\lfloor 8\{x\}\rfloor$ changes when $\{x\}=\frac{m}{n}$, where $m\in\{1,2,3,...,n-1\}$, $n\in\{2,4,6,8\}$. Using \href{https://artofproblemsolving.com/wiki/index.php?title=Euler\%27s\_Totient\_Function}{Euler's Totient Function}
 \[\sum\limits_{k=0}^4 \phi(2k)\]to obtain $12$ different values for $\{x\}=\frac{m}{n}$. (note that here Euler's Totient Function counts the number of $\{x\}=\frac{m}{n}$ where $m$, $n$ are relatively prime so that the values of $\{x\}$ won't overlap.). 

 Thus if $k$ can be expressed as $\lfloor 2x\rfloor+\lfloor 4x\rfloor+\lfloor 6x\rfloor+\lfloor 8x\rfloor$, then $k=20a+b$ for some non-negative integers $a$, $b$, where there are $12$ values for $b$. 

 Exclusively, there are $49$ values for $a$ in the range $0<k<1000$, or $49\cdot12=588$ ordered pairs $(a,b)$.

 If $a=0$, $b\neq0$, which includes $11$ ordered pairs.

 If $a=50$, $b=0$, which includes $1$ ordered pair.

 In total, there are $588+11+1=\boxed{600}$ values for $k$.

 





\textbf{Solution 5}

To simplify the question, let $y = 2x$. Then, the expression in the question becomes $\lfloor y \rfloor + \lfloor 2y \rfloor + \lfloor 3y \rfloor + \lfloor 4y \rfloor$. 

 Let $\{x\}$ represent the non-integer part of $x$ (For example, $\{2.8\} = 0.8$). Then, 

 
 \begin{align*} \lfloor y \rfloor + \lfloor 2y \rfloor + \lfloor 3y \rfloor + \lfloor 4y \rfloor &= y - \{y\} + 2y - \{2y\} + 3y - \{3y\} + 4y - \{4y\} \\ &= 10y - (\{y\} + \{2y\} + \{3y\} + \{4y\}) \\ &= 10(\lfloor y \rfloor + \{y\}) - (\{y\} + \{2y\} + \{3y\} + \{4y\}) \\ &= 10\lfloor y \rfloor + 10\{y\} - (\{y\} + \{2y\} + \{3y\} + \{4y\}) \\ &= 10\lfloor y \rfloor + 9\{y\} - (\{2y\} + \{3y\} + \{4y\}) \\ \end{align*}


 Since $\lfloor y \rfloor$ is always an integer, $10\lfloor y \rfloor$ will be a multiple of 10. Thus, we look for the range of the other part of the expression. We will be able to reach the same numbers when $y$ ranges from $0$ to $1$, because the curly brackets ($\{\}$) gets rid of any integer part. Let the combined integer part of $2y$, $3y$, and $4y$ be $k$ (In other words, $k = \lfloor 2y \rfloor + \lfloor 3y \rfloor + \lfloor 4y \rfloor$). Then, 

 
 \begin{align*} 9\{y\} - (\{2y\} + \{3y\} + \{4y\}) &= 9\{y\} - (2\{y\} + 3\{y\} + 4\{y\} - k) \\ &= 9\{y\} - (9\{y\} - k) \\ &= k \end{align*}


 The maximum value of $k$ will be when $y$ is slightly less than $1$, which means $k = 1 + 2 + 3 = 6$. As $y$ increases from $0$ to $1$, $k$ will increase whenever $2y$, $3y$, or $4y$ is an integer, which happens when $y$ hits one of the numbers in the set $\left\{\dfrac14, \dfrac13, \dfrac12, \dfrac23, \dfrac34 \right\}$. When $y$ reaches $\dfrac12$, both $2y$ and $4y$ will be an integer, so $k$ will increase by $2$. For all the other numbers in the set, $k$ increases by $1$ since only $1$ number in the set will be an integer. Thus, the possible values of $k$ are $\{0, 1, 2, 4, 5, 6\}$. 

 Finally, looking back at the original expression, we plug in $k$ to get a multiple of $10$ plus any number in $\{0, 1, 2, 4, 5, 6\}$. Thus, we hit all numbers ending in $\{0, 1, 2, 4, 5, 6\}$, of which there are $\boxed{600}$. 

 





\textbf{Solution 6}

Imagine that we gradually increase $x$ from $0$ to $1$. At the beginning, the value of our expression is $0$, at the end it is $2+4+6+8=20$. Note that every time $x=\frac{a}{b}$ for some positive integer $a$ and a positive multiple $b$ of either $2, 4, 6,$ or $8$. Thus, we have been able to express 12 of the integers from 1 through 20 when $0<x<1$, namely when 
 \[x = \frac{1}{2}, \frac{2}{2}=1, \frac{1}{4}, \frac{3}{4}, \frac{1}{6}, \frac{1}{3}, \frac{2}{3}, \frac{5}{6}, \frac{1}{8}, \frac{3}{8}, \frac{5}{8}, \frac{7}{8}\].

 
Notice that 
 \[\lfloor 2(x+n) \rfloor + \lfloor 4(x+n)\rfloor + \lfloor 6(x+n) \rfloor+ \lfloor8(x+n) \rfloor= \lfloor 2x \rfloor+ \lfloor 4x \rfloor+ \lfloor 6x \rfloor+ \lfloor 8x \rfloor+ 20n\]. This conceptually means that the number of integers that can be expressed in the range $(i, j)$ is the same as the number of integers that can be expressed in the range $(i+x, j+x)$. Thus, because we can express $12$ integers in the range $(1, 20)$, we can cover $12*50 = \boxed{600}$ from 1 to 1000.$-\text{thinker123}$

 





\textbf{Solution 7}

After observing, we can see that there are $6$ values of can be evaluated through the expression every $10$ numbers, so our answer is $6*100=600$ ~bluesoul

 




	\item (\textit{AIME I 1985}) An ellipse has foci at $(9, 20)$ and $(49, 55)$ in the $xy$-plane and is tangent to the $x$-axis. What is the length of its major axis?


 
 \textbf{Answer: }085



\textbf{Solution 1}

An ellipse is defined to be the \href{https://artofproblemsolving.com/wiki/index.php?title=Locus}{locus} of points $P$ such that the sum of the distances between $P$ and the two foci is constant.  Let $F_1 = (9, 20)$, $F_2 = (49, 55)$ and $X = (x, 0)$ be the point of tangency of the ellipse with the $x$-axis.  Then $X$ must be the point on the axis such that the sum $F_1X + F_2X$ is minimal.  (The last claim begs justification:  Let $F'_2$ be the reflection of $F_2$ across the $x$-axis. Let $Y$ be where the line through $F_1$ and $F’_2$ intersects the ellipse.  We will show that $X=Y$. Note that  $X F_2 = X F’_2$ since $X$ is on the $x$-axis.  Also, since the entire ellipse is on or above the $x$-axis and the line through $F_2$ and $F’_2$ is perpendicular to the $x$-axis, we must have $F_2 Y \leq F’_2 Y$ with equality if and only if $Y$ is on the $x$-axis.  Now, we have
 \[F_1 X + F'_2 X = F_1 X + F_2 X = F_1 Y + F_2 Y \leq F_1 Y + F’_2 Y\]But the right most sum is the straight-line distance from $F_1$ to $F’_2$ and the left is the distance of some path from $F_1$ to $F_2$., so this is only possible if we have equality and thus $X = Y$). Finding the optimal location for $X$ is a classic problem: for any path from $F_1$ to $X$ and then back to $F_2$, we can reflect (as above) the second leg of this path (from $X$ to $F_2$) across the $x$-axis.  Then our path connects $F_1$ to the reflection $F_2'$ of $F_2$ via some point on the $x$-axis, and this path will have shortest length exactly when our original path has shortest length.  This occurs exactly when we have a straight-line path, and by the above argument, this path passes through the point of tangency of the ellipse with the $x-$axis.  

 
                    \begin{figure}[H]
                        \centering
                        \includegraphics[width=0.5\textwidth]{Solution/AIME_1985_I_11_Solution_Figure_1.png}
                    \end{figure}
                    The sum of the two distances $F_1 X$ and $F_2X$ is therefore equal to the length of the segment $F_1F_2'$, which by the \href{https://artofproblemsolving.com/wiki/index.php?title=Distance\_formula}{distance formula} is just $d = \sqrt{(9 - 49)^2 + (20 + 55)^2} = \sqrt{40^2 + 75^2} = 5\sqrt{8^2 + 15^2} = 5\cdot 17 = 85$.  

 Finally, let $A$ and $B$ be the two endpoints of the major axis of the ellipse.  Then by symmetry $AF_1 = F_2B$ so $AB = AF_1 + F_1B = F_2B + F_1B = d$ (because $B$ is on the ellipse), so the answer is $\boxed{085}$.

 





\textbf{Solution 1 Simplified}

This assumes you know properties of ellipses.Like above, X is the point on the x-axis minimizing $XF_1+XF_2$.By reflection, this equals $F_1F_2'$...But this just equals the length of the major axis.Solving like above, we get $\boxed{085}$

 





\textbf{Solution 2 (Calculus)}

An ellipse is defined as the set of points where the sum of the distances from the foci to the point is fixed. The length of major axis is equal to the sum of these distances $(2a)$. Thus if we find the sum of the distances, we get the answer. Let k be this fixed sum; then we get, by the distance formula:

 $k = \sqrt{(x - 9)^2 + 20^2} + \sqrt{(x - 49)^2 + 55^2}$

 This is the equation of the ellipse expressed in terms of $x$. The line tangent to the ellipse at the given point $P(x, 0)$ will thus have slope $0$. Taking the derivative gives us the slope of this line. To simplify, let $f(x) = (x - 9)^2 + 20^2$ and $g(x) = (x - 49)^2 + 55^2$. Then we get:

 $0 = \frac{f^\prime(x)}{2\sqrt{f(x)}} + \frac{g^\prime(x)}{2\sqrt{g(x)}}$

 Next, we multiply by the conjugate to remove square roots. We next move the resulting $a^2 - b^2$ form expression into form $a^2 = b^2$.

 $\frac{(f^\prime(x))^2}{4\cdot f(x)} = \frac{(g^\prime(x))^2}{4\cdot g(x)}$ 

 We know $f^\prime(x) = 2x - 18$ and $g^\prime(x) = 2x - 98$. Simplifying yields:

 $\frac{(x - 9)^2}{(x - 9)^2 + 20^2} = \frac{(x - 49)^2}{(x - 49)^2 + 55^2}$

 To further simplify, let $a = (x - 9)^2$ and $b = (x - 49)^2$. This means $\frac{a}{a + 400} = \frac{b}{b + 3025}$. Solving yields that $16b = 121a$.Substituting back $a$ and $b$ yields:

 $16 \cdot (x - 49)^2 = 121 \cdot (x - 9)^2$.

 Solving for $x$ yields $x = \frac{59}{3}$. Substituting back into our original distance formula, solving for $k$ yields $\boxed{085}$.

 \href{https://mathproblemsolvingskills.wordpress.com/}{MathProblemSolvingSkills.com}

 


 




	\item (\textit{AIME I 1985}) Let $A$, $B$, $C$ and $D$ be the vertices of a regular tetrahedron, each of whose edges measures $1$ meter. A bug, starting from vertex $A$, observes the following rule: at each vertex it chooses one of the three edges meeting at that vertex, each edge being equally likely to be chosen, and crawls along that edge to the vertex at its opposite end. Let $p = \frac{n}{729}$ be the probability that the bug is at vertex $A$ when it has crawled exactly $7$ meters. Find the value of $n$.


 
 \textbf{Answer: }182



\textbf{Solution 1 (Single Variable Recursion)}

For all nonnegative integers $k,$ let $P(k)$ be the probability that the bug is at vertex $A$ when it has crawled exactly $k$ meters. We wish to find $p=P(7).$

 Clearly, we have $P(0)=1.$ For all $k\geq1,$ note that after $k-1$ crawls:

 <ol style="margin-left: 1.5em;"> \item The probability that the bug is at vertex $A$ is $P(k-1),$ and the probability that it crawls to vertex $A$ on the next move is $0.$ 

 \item The probability that the bug is not at vertex $A$ is $1-P(k-1),$ and the probability that it crawls to vertex $A$ on the next move is $\frac13.$

 </ol>Together, the recursive formula for $P(k)$ is
 \[P(k) = \begin{cases} 1 & \mathrm{if} \ k=0 \\ \frac13(1-P(k-1)) & \mathrm{if} \ k\geq1 \end{cases}.\]Two solutions follow from here:

 





\textbf{Solution 1.1 (Recursive Formula)}

We evaluate $P(7)$ recursively:
 \begin{alignat*}{6} P(0)&=1, \\ P(1)&=\frac13(1-P(0))&&=0, \\ P(2)&=\frac13(1-P(1))&&=\frac13, \\ P(3)&=\frac13(1-P(2))&&=\frac29, \\ P(4)&=\frac13(1-P(3))&&=\frac{7}{27}, \\ P(5)&=\frac13(1-P(4))&&=\frac{20}{81}, \\ P(6)&=\frac13(1-P(5))&&=\frac{61}{243},\\ P(7)&=\frac13(1-P(6))&&=\frac{182}{729}. \end{alignat*}Therefore, the answer is $n=\boxed{182}.$

 





\textbf{Solution 1.2 (Explicit Formula)}

Let $P(k)=Q(k)+c$ for some function $Q(k)$ and constant $c.$ For all $k\geq1,$ the recursive formula for $P(k)$ becomes 
 \[Q(k)+c=\frac13(1-(Q(k-1)+c))=\frac13-\frac13Q(k-1)-\frac13c.\] Solving for $Q(k),$ we get 
 \[Q(k)=\frac13-\frac13Q(k-1)-\frac43c.\]For simplicity purposes, we set $c=\frac14,$ which gives 
 \[Q(k)=-\frac13Q(k-1).\]Clearly, $Q(0),Q(1),Q(2),\ldots$ is a geometric sequence with the common ratio $\frac{Q(k)}{Q(k-1)}=-\frac13.$ Substituting $P(0)=1$ and $c=\frac14$ into $P(0)=Q(0)+c$ produces $Q(0)=\frac34,$ the first term of the geometric sequence. 

 For all nonnegative integers $k,$ the explicit formula for $Q(k)$ is 
 \[Q(k)=\frac34\left(-\frac13\right)^k,\] and the explicit formula for $P(k)$ is 
 \[P(k)=\frac34\left(-\frac13\right)^k+\frac14.\] Finally, the requested probability is $p=P(7)=\frac{182}{729},$ from which $n=\boxed{182}.$

 





\textbf{Solution 2 (Multivariable Recursion by Algebra)}

<u>\textbf{Denominator}</u>

 There are $3^7$ ways for the bug to make $7$ independent crawls without restrictions.

 <u>\textbf{Numerator}</u>

 Let $V_k$ denote the number of ways for the bug to crawl exactly $k$ meters starting from vertex $V$ and ending at vertex $A,$ where $V\in\{A,B,C,D\}$ and $k$ is a positive integer. We wish to find $A_7.$

 Since the bug must crawl to vertex $B,C,$ or $D$ on the first move, we have
 \begin{align*} A_7&=B_6+C_6+D_6 \\ &=(A_5+C_5+D_5)+(A_5+B_5+D_5)+(A_5+B_5+C_5) \\ &=A_5+2(A_5+B_5+C_5+D_5) \\ &=A_5+2S_5, \end{align*}
where $S_k=A_k+B_k+C_k+D_k.$

 More generally, we get 
 \[A_{k+2}=A_k+2S_k. \qquad\qquad (\spadesuit)\]For $S_k,$ note that

 
            \begin{enumerate}
                 \item Base Case:
 \begin{align*} S_1&=A_1+B_1+C_1+D_1 \\ &=0+1+1+1 \\ &=3. \end{align*}
 \item Recursive Case:
 \begin{align*} S_{k+1}&=A_{k+1}+B_{k+1}+C_{k+1}+D_{k+1} \\ &=(B_k+C_k+D_k)+(A_k+C_k+D_k)+(A_k+B_k+D_k)+(A_k+B_k+C_k) \\ &=3(A_k+B_k+C_k+D_k) \\ &=3S_k. \end{align*}

            \end{enumerate}
            Clearly, $S_1,S_2,S_3,\ldots$ is a geometric sequence with the first term $S_1=3$ and the common ratio $\frac{S_{k+1}}{S_k}=3.$ Therefore, its explicit formula is 
 \[S_k=3^k. \qquad\qquad (\clubsuit)\]Recall that $A_1=0.$ By $(\spadesuit)$ and $(\clubsuit),$ we rewrite $A_7$ recursively:
 \begin{align*} A_7 &= A_5+2S_5 \\ &=\left(A_3+2S_3\right)+2S_5 \\ &=\left(\left(A_1+2S_1\right)+2S_3\right)+2S_5 \\ &=2S_1+2S_3+2S_5 \\ &=2(3)+2\left(3^3\right)+2\left(3^5\right). \end{align*}
<u>\textbf{Answer}</u>

 The requested probability is 
 \[p=\frac{A_7}{3^7}=\frac{2(3)+2\left(3^3\right)+2\left(3^5\right)}{3^7}=\frac{2(1)+2\left(3^2\right)+2\left(3^4\right)}{3^6}=\frac{182}{729},\] from which $n=\boxed{182}.$

 





\textbf{Solution 3 (Multivariable Recursion by Table)}

Define notation $V_k$ as Solution 2 does.

 In fact, we can generalize the following relationships for all <i>\textbf{nonnegative}</i> integers $k:$
 \begin{align*} A_0&=1, \\ B_0&=0, \\ C_0&=0, \\ D_0&=0, \\ A_{k+1}&=B_k+C_k+D_k, \\ B_{k+1}&=A_k+C_k+D_k, \\ C_{k+1}&=A_k+B_k+D_k, \\ D_{k+1}&=A_k+B_k+C_k. \\ \end{align*}
Using these equations, we recursively fill out the table below:
 \[\begin{array}{c||c|c|c|c|c|c|c|c}    \hspace{7mm}&\hspace{6mm}&\hspace{6mm}&\hspace{6mm}&\hspace{6mm}&\hspace{6mm}&\hspace{6mm}&& \\ [-2.5ex] \boldsymbol{k} & \boldsymbol{0} & \boldsymbol{1} & \boldsymbol{2} & \boldsymbol{3} & \boldsymbol{4} & \boldsymbol{5} & \boldsymbol{6} & \boldsymbol{7} \\  \hline \hline &&&&&&&& \\ [-2.25ex] \boldsymbol{A_k} &1&0&3&6&21&60&183&546 \\ \hline   &&&&&&&& \\ [-2.25ex] \boldsymbol{B_k} &0&1&2&7&20&61&182&547 \\ \hline  &&&&&&&& \\ [-2.25ex] \boldsymbol{C_k} &0&1&2&7&20&61&182&547 \\ \hline    &&&&&&&& \\ [-2.25ex] \boldsymbol{D_k} &0&1&2&7&20&61&182&547 \\ \end{array}\]Note that the paths from $V$ to $A$ and the paths from $A$ to $V$ have one-to-one correspondence. So, we must get 
 \[A_k+B_k+C_k+D_k=3^k\] for all values of $k.$

 The requested probability is 
 \[p=\frac{A_7}{3^7}=\frac{546}{2187}=\frac{182}{729},\] from which $n=\boxed{182}.$

 





\textbf{Solution 4 (Single Variable Version of Solution 2)}

Let $a_n$ denotes the number of ways that the bug arrives at $A$ after crawling $n$ meters, then we have $a_1=0$.

 Notice that there is respectively $1$ way to arrive at $A$ for each of the different routes after the previous $n-1$ crawls, excluding the possibility that the bug ends up at $A$ after the $(n-1)$th crawl (as it will be forced to move somewhere else.). Thus, we get the recurrence relation 
 \[a_n=3^{n-1}-a_{n-1}.\]Quick calculations yield
 \begin{align*} a_7&=3^6-a_6\\ &=3^6-\left(3^5-3^4+3^3-3^2+3-a_1\right)\\ &=546. \end{align*}
Thus, $p=\frac{546}{3^7}=\frac{182}{729}\Longrightarrow n=\boxed{182}$.

 





\textbf{Solution 5 (Dynamic Programming)}

Let $A(n)$ be the probability the bug lands on vertex $A$ after crawling $n$ meters, $B(n)$ be the probability the bug lands on vertex $B$ after crawling $n$ meters, and etc.

 Note that $A(1)=0$ and $B(1)=C(1)=D(1)=\frac13.$ For $n\geq2,$ the probability that the bug land on each vertex after $n$ meters is $\frac13$ the sum of the probability the bug land on other vertices after crawling $n-1$ meters. So, we have
 \begin{align*} A(n) &= \frac13 \cdot [B(n-1) + C(n-1) + D(n-1)], \\ B(n) &= \frac13 \cdot [A(n-1) + C(n-1) + D(n-1)], \\ C(n) &= \frac13 \cdot [A(n-1) + B(n-1) + D(n-1)], \\ D(n) &= \frac13 \cdot [A(n-1) + B(n-1) + C(n-1)]. \end{align*}
It follows that $A(n) = B(n-1) = C(n-1) = D(n-1).$

 We construct the following table:
 \[\begin{array}{c|cccc}  &  &  &  & \\ [-2ex] n & A(n) & B(n) & C(n) & D(n) \\ [1ex] \hline  &  &  &  & \\ [-1ex] 1 & 0 & \frac13 & \frac13 & \frac13 \\  &  &  &  & \\ 2 & \frac13 & \frac29 & \frac29 & \frac29 \\  &  &  &  & \\ 3 & \frac29 & \frac{7}{27} & \frac{7}{27} & \frac{7}{27} \\  &  &  &  & \\ 4 & \frac{7}{27} & \frac{20}{81} & \frac{20}{81} & \frac{20}{81}\\  &  &  &  & \\ 5 & \frac{20}{81} & \frac{61}{243} & \frac{61}{243} & \frac{61}{243} \\  &  &  &  & \\ 6 & \frac{61}{243} & \frac{182}{729} & \frac{182}{729} & \frac{182}{729} \\  &  &  &  & \\ 7 & \frac{182}{729} & \frac{547}{2187} & \frac{547}{2187} & \frac{547}{2187} \\ [1ex] \end{array}\]Therefore, the answer is $n = \boxed{182}.$

 





\textbf{Solution 6 (Generating Functions and Roots of Unity Filter / Casework)}

The generating function for a problem with this general form ($4$ states, $n$ steps) is $(x+x^2+x^3)^n$, so the generating function of interest for this problem is $(x+x^2+x^3)^7$. Our goal is to find the coefficients of every $x^{4n}$ and add them up before dividing by $3^7$. Here we have two ways to proceed:

 \textbf{1. Roots of Unity Filter} Let $\omega = e^{i\pi / 2}$. We have that if $G(x) = (x+x^2+x^3)^7$, then 
 \[\frac{G(1) + G(\omega) + G(\omega^2) + G(\omega^3)}{4} = \frac{2187-1-1-1}{4} = 546.\] From here, the desired probability is $\frac{546}{2187} = \frac{182}{729}$. Therefore, the answer is $n=\boxed{182}$.

 \textbf{2. Casework} We can factor $(x+x^2+x^3)^7$ as $x^7(1+x+x^2)^7.$ The $x^{4n}$ coefficients of $(x+x^2+x^3)^7$ will be the same as the $x^{4n+1}$ coefficients of $(1+x+x^2)^7.$ The possible values for $4n+1$ would then be $1,$ $5,$ $9,$ and $13.$ Because $1+13=5+9=14,$ the coefficients of $x^1$ and $x^{13}$ are equal and so are the coefficients of $x^5$ and $x^9.$ To make an $x$ term, we need $6$ "$1$" terms and one "$x$" term multiplied together. There are $7$ ways to arrange these numbers. The coefficient of the $x^5$ term will be the sum of the number of the possible arrangements for these numbers:
 \begin{align*} 0000122&=105\text{ arrangements}, \\ 0001112&=140\text{ arrangements}, \\ 0011111&=21\text{ arrangements}. \end{align*}
Thus, the coefficient of the $x^5$ term is $105+140+21=266.$ From here, the desired probability is $\frac{2(7+266)}{2187}=\frac{546}{2187}=\frac{182}{729}.$ Thus, our answer is $\boxed{182}.$

 





\textbf{Solution 7 (Partitions)}

We can find the number of different times the bug reaches vertex $A$ before the $7$th move, and use these smaller cycles to calculate the number of different ways the bug can end up back at $A.$

 Define $f(x)$ to be the number of paths of length $x$ which start and end at $A$ but do not pass through $A$ otherwise. Obviously $f(1) = 0.$ In general for $f(x),$ the bug has three initial edges to pick from. From there, since the bug cannot return to $A$ by definition, the bug has exactly two choices. This continues from the $2$nd move up to the $(x-1)$th move. The last move must be a return to $A,$ so this move is determined. So $f(x) = 2^{x-2}3.$

 Now we need to find the number of cycles by which the bug can reach $A$ at the end. Since $f(1) = 0,$ we know that $f(6)$ cannot be used, as on the $7$th move the bug cannot move from $A$ to $A.$ So we need to find the number of \href{https://artofproblemsolving.com/wiki/index.php?title=Partition}{partitions} of $7$ using only $2,3,4,5,$ and $7.$ These are $f(2)f(2)f(3),f(2)f(5),f(3)f(4),$ and $f(7).$ We can calculate these and sum them up using our formula. Also, order matters, so we need to find the number of ways to arrange each partition:
 \[{3\choose1}f(2)f(2)f(3) + {2\choose1}f(2)f(5) + {2\choose1}f(3)f(4) + f(7) = 3(3)(3)(2\cdot3) + 2(3)(2^33) + 2(2\cdot3)(2^23) + (2^53) = 546.\]Finally, this is a probability question, so we divide by $3^7,$ or $\frac{546}{3^7} = \frac{182}{3^6}.$ The answer is $n=\boxed{182}.$

 





\textbf{Solution 8 (Sequences)}

Note that this problem is basically equivalent to the following:How many distinct sequences of $8$ integers $a_1, a_2, a_3, \ldots, a_8$ are there such that $a_1 = a_8 = 1,$ $a_i \in \{1, 2, 3, 4\}$ for all $2 \leq i\leq 8,$ and $a_i \neq a_{i+1}$ for all $1 \leq i \leq 7$? 

 Now consider the $8$ integers modulo $4.$ Let $b_1, b_2, b_3, \ldots, b_7$ be a new sequence of integers such that $b_i = a_{i+1} - a_i \mod 4$ for all $1 \leq i \leq 7.$

 Note that the condition is equivalent to that $b_i \in \{1, 2, 3\}$ for all $1 \leq i \leq 7,$ and since $a_1 \mod 4 = a_8 \mod 4,$ $b_1 + b_2 + \cdots + b_7$ must be a multiple of $4.$

 Thus, our desired number of paths is equivalent to the number of ordered septuples of positive integers $(b_1, b_2, \ldots, b_7)$ such that $b_i \in \{1, 2, 3\}$ for all $1 \leq i \leq 7$ and $b_1 + b_2 + \cdots + b_7$ is congruent to $0 \mod 4.$

 We will now proceed with casework on the number of $2$s in the septuple.

 One $2$: There are $\dbinom{7}{1} = 7$ ways to arrange the $2$, and $\dbinom{6}{0} + \dbinom{6}{2} + \dbinom{6}{4} + \dbinom{6}{6} = 2^5$ valid ways (the proof of this combinatorial identity is left as an exercise to the reader) to arrange the $1$s and $3$s.

 Three $2$s: There are $\dbinom{7}{3} = 35$ ways to arrange the $2$s, and $\dbinom{4}{1} + \dbinom{4}{3} = 2^3$ valid ways to arrange the $1$s and $3$s.

 Five $2$s: There are $\dbinom{7}{5} = 21$ ways to arrange the $2$s, and $\dbinom{2}{0} + \dbinom{2}{2} = 2$ valid ways to arrange the $1$s and $3$s.

 Adding up our cases, we obtain $7 \cdot 32 + 35 \cdot 8 + 21 \cdot 2 = 546$ valid sequences.Dividing by the total number of paths without restrictions, $3^7 = 2187,$ our desired probability is $\frac{546}{2187} = \frac{182}{729},$ giving an answer of $\boxed{182}.$

 





\textbf{Solution 9}

We instead find the probability that the bug is NOT at vertex $A$ after crawling $n$ meters (equivalent to moving $n$ times). Call $A_n$ the probability that the bug IS at vertex $A$ after $n$ moves; call $O_n$ the probability that the bug is on some other vertex. We have the following recurrence relations. 
 \[A_n = \frac{1}{3}O_{n-1}\] 
 \[O_n = A_{n-1} + \frac{2}{3}O_{n-1}\] However, we can calculate $A_{n-1}$ in terms of $O_n$; take $n = k-1$, and we have 
 \[A_{k-1} = \frac{1}{3}O_{k-2}\]. Substituting this into our equation for $O$, we have that 
 \[O_n = \frac{1}{3}O_{n-2} + \frac{2}{3}O_{n-1}\]. We also have the conditions that $O_0 = 0$ (the bug will not be at vertex other than $A$ on the "0th" move) and $O_1 = 1$ (the bug will be at a vertex other than $A$ after the $1st$ move). Iteratively calculating $O_7$, we find that it is equal to $\frac{547}{729}$ (I will not do this calculation here; you can do it manually if you wish to check). However, this is the probability that the ant is NOT at vertex $A$ after $7$ moves; then its complement, $\frac{182}{729}$ is the probability that the ant IS at vertex $A$ after $7$ moves, so our answer is $\boxed{182}$.

 





\textbf{Solution 10 (Linear Algebra)}

Think of the problem as a state machine with a transition matrix. State 1 is if the bug is at A, State 2 is if the bug is not at A. In vector form we can represent state 1 as $\begin{bmatrix} 1 \\ 0 \end{bmatrix}$ and state 2 as $\begin{bmatrix} 0 \\ 1 \end{bmatrix}$. Our transition matrix $T =\begin{pmatrix} 0 & \frac{1}{3} \\ 1 & \frac{2}{3} \end{pmatrix}$. Thus the probability of being in each state after k moves is $T^{k}\begin{bmatrix} 1 \\ 0 \end{bmatrix}=\begin{pmatrix} 0 & \frac{1}{3} \\ 1 & \frac{2}{3} \end{pmatrix}^{k}\begin{bmatrix} 1 \\ 0 \end{bmatrix}$. Those familiar with linear algebra will recognize the need to diagonalize the transition matrix through eigendecomposition. In short, the eigenvalues of the transition matrix are $-\frac{1}{3}$ and 1, with corresponding eigenvector basis of $\begin{pmatrix} 1 \\ -1 \end{pmatrix},  \begin{pmatrix} \frac{1}{3} \\ 1 \end{pmatrix}$. Thus $T = Q \Lambda Q^{-1} = \begin{pmatrix} 1 & \frac{1}{3} \\ -1 & 1 \end{pmatrix}\begin{pmatrix} -\frac{1}{3} & 0\\ 0 & 1 \end{pmatrix}\begin{pmatrix} 1 & \frac{1}{3}\\ -1 & 1 \end{pmatrix}^{-1} = \begin{pmatrix} 1 & \frac{1}{3} \\ -1 & 1 \end{pmatrix}\begin{pmatrix} -\frac{1}{3} & 0\\ 0 & 1 \end{pmatrix}\begin{pmatrix} \frac{3}{4} & -\frac{1}{4}\\ \frac{3}{4} & \frac{3}{4} \end{pmatrix}$

 Thus $T^{k}\begin{bmatrix} 1 \\ 0 \end{bmatrix}= \begin{pmatrix} 1 & \frac{1}{3} \\ -1 & 1 \end{pmatrix}\begin{pmatrix} (-\frac{1}{3})^{k} & 0\\ 0 & 1^{k} \end{pmatrix}\begin{pmatrix} \frac{3}{4} & -\frac{1}{4}\\ \frac{3}{4} & \frac{3}{4} \end{pmatrix}\begin{bmatrix} 1 \\ 0 \end{bmatrix} = \begin{pmatrix} \frac{3}{4}(-\frac{1}{3})^{k}+\frac{1}{4} & -\frac{1}{4}(-\frac{1}{3})^{k}+\frac{1}{4}\\ -\frac{3}{4}(-\frac{1}{3})^{k}+\frac{3}{4} &  \frac{1}{4}(-\frac{1}{3})^{k}+\frac{3}{4} \end{pmatrix}\begin{bmatrix} 1 \\ 0 \end{bmatrix} = \begin{bmatrix} \frac{3}{4}(-\frac{1}{3})^{k}+\frac{1}{4} \\ -\frac{3}{4}(-\frac{1}{3})^{k}+\frac{3}{4} \end{bmatrix}$. 

 The probability we seek is the top entry of this vector for k = 7, or $\frac{182}{729}$.

 





\textbf{Solution 11 (Regular Casework)}

Perhaps you are not very good at recursion. Let $A$ denote the vertex $A,$ and let $N$ denote a vertex other than $A.$ Now the moves of the bug must look something like this:
 \[N - - - - - A.\]

 Note that the bug cannot travel to $A$ twice in a row, meaning it can't have an $A$ following an $A,$ so there must be at least one $N$ in between. Therefore, the bug can only move to $A$ a maximum of $3$ times (including the last one). Therefore, we will do casework based on how many times the bug moves to vertex $A.$$\newline$$\newline$Case $1: \text{3 As}$

 There are three ways for the bug to move in this case:
 \[\text{N A N A N N A},\]
 \[\text{N A N N A N A},\]
 \[\text{N N A N A N A}.\]$\newline$

 Let $P(XY)$ denote the probability that the bug moves from vertex $X$ to vertex $Y.$$\newline$$P(NA) = \frac{1}{3},$ because from a vertex $N$ there are three ways to go, with only one of them being $A.$$\newline$$P(AN) = 1,$ self explanatory.$\newline$$P(NN) = \frac{2}{3},$ also self explanatory.$\newline$Now, using this (by multiplying the probabilities for each pair of consecutive letters), we find that each configuration has the same probability: $(P(NN))(P(NA))^3 = \frac{2}{81}.$ Since each of these three configurations has this probability, we only need to multiply by $3,$ giving $\frac{2(3)}{81} = \frac{2}{27}$ so far.$\newline$$\newline$Case $2: \text{2 As}$

 This is a little simpler.The configuration:
 \[N N N N A N A,\]where $A$ can occupy any space from the second to the fifth. Note that each occupation gives the same probability once again, so we have $(5-(2-1))\times (P(NN))^3(P(NA))^2 = 4\times \left(\frac{2}{3}\right)^3\left(\frac{1}{3}\right)^2 = \frac{32}{243}.$$\newline$$\newline$Case $3:\text{1 A}$Configuration:

 
 \[N N N N N N A\]

 The probability of this is simply $(P(NN))^5(P(NA)) = \frac{2^5}{3^6} = \frac{32}{729}.$

 Adding this all up, we have 
 \[\frac{2}{27} + \frac{32}{243} + \frac{32}{729} = \frac{54+96+32}{729} = \frac{182}{729} \Longrightarrow \boxed{182}.\]

 \textbf{Remark} Here is a similar problem from another AIME test: \href{https://artofproblemsolving.com/wiki/index.php?title=2003\_AIME\_II\_Problems/Problem\_13}{2003 AIME II Problem 13}, in which we have an equilateral triangle instead.

 There is another similar problem from the AMC8: \href{https://artofproblemsolving.com/wiki/index.php?title=2022\_AMC\_8\_Problems/Problem\_25}{2022 AMC 8 Problems/Problem 25}, where we have the same question, just with less steps

 




	\item (\textit{AIME I 1985}) The numbers in the \href{https://artofproblemsolving.com/wiki/index.php?title=Sequence}{sequence} $101$, $104$, $109$, $116$,$\ldots$ are of the form $a_n=100+n^2$, where $n=1,2,3,\ldots$. For each $n$, let $d_n$ be the greatest common divisor of $a_n$ and $a_{n+1}$. Find the maximum value of $d_n$ as $n$ ranges through the \href{https://artofproblemsolving.com/wiki/index.php?title=Positive\_integer}{positive integers}.


 
 \textbf{Answer: }401



\textbf{Solution 1}

If $(x,y)$ denotes the \href{https://artofproblemsolving.com/wiki/index.php?title=Greatest\_common\_divisor}{greatest common divisor} of $x$ and $y$, then we have $d_n=(a_n,a_{n+1})=(100+n^2,100+n^2+2n+1)$. Now assuming that $d_n$ \href{https://artofproblemsolving.com/wiki/index.php?title=Divisor}{ divides} $100+n^2$, it must divide $2n+1$ if it is going to divide the entire \href{https://artofproblemsolving.com/wiki/index.php?title=Expression}{expression} $100+n^2+2n+1$.

 Thus the \href{https://artofproblemsolving.com/wiki/index.php?title=Equation}{equation} turns into $d_n=(100+n^2,2n+1)$. Now note that since $2n+1$ is \href{https://artofproblemsolving.com/wiki/index.php?title=Odd\_integer}{ odd} for \href{https://artofproblemsolving.com/wiki/index.php?title=Integer}{ integral} $n$, we can multiply the left integer, $100+n^2$, by a power of two without affecting the greatest common divisor. Since the $n^2$ term is quite restrictive, let's multiply by $4$ so that we can get a $(2n+1)^2$ in there.

 So $d_n=(4n^2+400,2n+1)=((2n+1)^2-4n+399,2n+1)=(-4n+399,2n+1)$. It simplified the way we wanted it to!Now using similar techniques we can write $d_n=(-2(2n+1)+401,2n+1)=(401,2n+1)$.  Thus $d_n$ must divide $\boxed{401}$ for every single $n$.  This means the largest possible value for $d_n$ is $401$, and we see that it can be achieved when $n = 200$.

 





\textbf{Solution 2}

We know that $a_n = 100+n^2$ and $a_{n+1} = 100+(n+1)^2 = 100+ n^2+2n+1$. Since we want to find the GCD of $a_n$ and $a_{n+1}$, we can use the \href{https://artofproblemsolving.com/wiki/index.php?title=Euclidean\_algorithm}{Euclidean algorithm}:

 $a_{n+1}-a_n = 2n+1$

 Now, the question is to find the GCD of $2n+1$ and $100+n^2$.We subtract $2n+1$ $100$ times from $100+n^2$.
 \[(100+n^2)-100(2n+1)\]
 \[=n^2+100-200n-100\]This leaves us with 
 \[n^2-200n.\]Factoring, we get 
 \[n(n-200)\]Because $n$ and $2n+1$ will be coprime, the only thing stopping the GCD from being $1$ is $n-200.$We want this to equal 0, because that will maximize our GCD. Solving for $n$ gives us $n=200$. The last remainder is 0, thus $200*2+1 = \boxed{401}$ is our GCD.

 





\textbf{Solution 3}

If Solution 2 is not entirely obvious, our answer is the max possible range of $\frac{x(x-200)}{2x+1}$. Using the Euclidean Algorithm on $x$ and $2x+1$ yields that they are relatively prime. Thus, the only way the GCD will not be 1 is if the$x-200$ term share factors with the $2x+1$. Using the Euclidean Algorithm, $\gcd(x-200,2x+1)=\gcd(x-200,2x+1-2(x-200))=\gcd(x-200,401)$. Thus, the max GCD is 401.

 





\textbf{Solution 4}

We can just plug in Euclidean algorithm, to go from $\gcd(n^2 + 100, n^2 + 2n + 101)$ to $\gcd(n^2 + 100, 2n + 1)$ to $\gcd(n^2 + 100 - 100(2n + 1), 2n + 1)$ to get $\gcd(n^2 - 200n, 2n + 1)$. Now we know that no matter what, $n$ is relatively prime to $2n + 1$. Therefore the equation can be simplified to: $\gcd(n - 200, 2n + 1)$. Subtracting $2n - 400$ from $2n + 1$ results in $\gcd(n - 200,401)$. The greatest possible value of this is $\boxed{401}$, and happens when $n \equiv 200 \pmod{401}$.

 





\textbf{Solution 5}

As clearly shown in the above solutions, we want to maximize $(100+n^2, 2n+1).$ Then the maximum of the GCD will be achieved with integers $a,b,c$ such that $a(100+n^2) + b(2n+1)=c$ by Bezout's identity. Note for the LHS to be constant, $b$ must be a linear function of $n,$ so $b=px+q.$ However, there cannot be a linear $n$ term in $b(2n+1),$ hence $b=2n-1$ by difference of squares. Changing $b$ to $-2n+1,$ we get $100a+an^2-4n^2+1=c,$ so $a=4,$ and the answer is $c=\boxed{401}.$ 

 




	\item (\textit{AIME I 1985}) In a tournament each player played exactly one game against each of the other players. In each game the winner was awarded 1 point, the loser got 0 points, and each of the two players earned $\frac{1}{2}$ point if the game was a tie. After the completion of the tournament, it was found that exactly half of the points earned by each player were earned in games against the ten players with the least number of points. (In particular, each of the ten lowest scoring players earned half of her/his points against the other nine of the ten). What was the total number of players in the tournament?


 
 \textbf{Answer: }025



\textbf{Solution 1}

Let us suppose for convenience that there were $n + 10$ players overall.  Among the $n$ players not in the weakest 10 there were $n \choose 2$ games played and thus $n \choose 2$ points earned.  By the givens, this means that these $n$ players also earned $n \choose 2$ points against our weakest 10.  Now, the 10 weakest players playing amongst themselves played ${10 \choose 2} = 45$ games and so earned 45 points playing each other.  Then they also earned 45 points playing against the stronger $n$ players.  Since every point earned falls into one of these categories, It follows that the total number of points earned was $2{n \choose 2} + 90 = n^2 - n + 90$.  However, there was one point earned per game, and there were a total of ${n + 10 \choose 2} = \frac{(n + 10)(n + 9)}{2}$ games played and thus $\frac{(n + 10)(n + 9)}{2}$ points earned.  So we have $n^2 -n + 90 = \frac{(n + 10)(n + 9)}{2}$ so $2n^2 - 2n + 180 = n^2 + 19n + 90$ and $n^2 -21n + 90 = 0$ and $n = 6$ or $n = 15$.  Now, note that the top $n$ players got $n(n - 1)$ points in total (by our previous calculation) for an average of $n - 1$, while the bottom 10 got 90 points total, for an average of 9.  Thus we must have $n > 10$, so $n = 15$ and the answer is $15 + 10 = \boxed{25}$.

 





\textbf{Solution 2}

Suppose that there are $n$ players participating in the tournament. We break this up into a group of the weakest ten, and the other $n-10$ people. Note that the $10$ players who played each other generated a total of $\dbinom{10}{2} = 45$ points playing each other. Thus, they earned $45$ playing the $n-10$ other people. Thus, the $n-10$ people earned a total of $10(n-10)-45 = 10n-145$ points playing vs. this group of 10 people, and also earned a total of $10n-145$ playing against themselves. Since each match gives a total of one point, we must have that $\dbinom{n-10}{2}=10n-145$. Expanding and simplifying gives us $n^2-41n+400=0$. Thus, $n=16$ or $n=25$. Note however that if $n=16$, then the strongest $16$ people get a total of $16*10-145=15$ playing against the weakest $10$ who gained $45$ points vs them, which is a contradiction since it must be larger. Thus, $n=\boxed{25}$.

 Solution by GameMaster402

 





\textbf{Solution 3}

Note that the total number of points accumulated must sum to ${p \choose 2} = \frac{p(p-1)}{2}$. Say the number of people is $n$. Consider the number of points gained when the 10 lowest scoring people play each other. The problem tells us that each of these 10 people must earn exactly half of the total number of points they will earn during the whole game. This implies that this group of 10 people must accumulate half their total combined points after they (the 10 people) all play each other, meaning they must earn the other half of their points by playing the $n-10$ stronger players. The problem also tells us that the $n-10$ people who aren't part of the losers group will earn half of their points by playing the $10$ losers.  Since the $n-10$ group and $10$ losers will earn half their points by playing each other, the sum of the number of points that they gain playing each other must then be half of the total amount of points earned by everyone in the game. Therefore, $\frac{p(p-1)}{4} = 10(p-10)$. This equation is the same as above, and by the same logic, the answer is $n=\boxed{25}$.

 





\textbf{Solution 4 (Expected Value and simple system of equations)}

Against any given person player \( p_1 \) has a \( 1/3 \) chance of obtaining \( 1 \) point, \( 1/3 \) chance of obtaining \( 0 \) points, and \( 1/3 \) chance of obtaining \( 1/2 \) a point. Thus, the expected value of the number of points earned in one game is just \( \frac{1}{3} \cdot 1 + \frac{1}{3} \cdot 0 + \frac{1}{3} \cdot \frac{1}{2} = \frac{1}{2} \). 

 Say there are \( n+10 \) people. Then, the number of games is determined by \( \binom{n+10}{2} \). The expected number of points for player \( p_1 \) is determined as 
 \[p_1 = \frac{1}{2} \binom{n+10}{2}\]

 Now, let the \( n+10 \) people be \( p_1, p_2, p_3, \dots, p_{n+10} \). The expected number of total points \( P \) is determined as the sum of all \( p_1, p_2, p_3, \dots, p_{n+10} \), or \( \left( \frac{1}{2} \binom{n+10}{2} \right)(n+10) \). 

 Now we use the other information given. We have that half the total number of points \( P \) came from the \( 10 \) people with the least points. WLOG say these people were \( p_1, p_2, p_3, \dots, p_{10} \). Thus we have
 \[\frac{1}{2} P = 10 \left( \frac{1}{2} \binom{10}{2} \right) + n \left( \frac{1}{2} \binom{n}{2} \right) + 5n\]

 If you don't understand this, let's break it down LHS to RHS. On the LHS, it is quite obvious. Half the total points obtain, which is $\frac{1}{2} P$. Now, first term on the RHS. Note what the MAA says in the brackets. The property applies for the 10 lowest. Thus, for each of \( p_1, p_2, p_3, \dots, p_{10} \), we must consider their individual chances of winning. Between their respective 10 player bracket, $\binom{10}{2}$ games occur. Then, the expected number of points gained per player is just $\frac{1}{2}$. Combining them together, we obtain the first term. Now, the middle term. We now consider the $n$ people going against each other. This time, we consider each of \( p_{11}, p_{12}, p_{13}, \dots, p_{n+10} \), which is $n$ total. Then, there are now $\binom{n}{2}$ games in the bracket, and each player has a $\frac{1}{2}$ chance of winning. Combining all together gives us the middle term. The last term involves each of the $n$ people going against the 10 lowest people. Then, $p_{11}$ goes against each of \( p_1, p_2, p_3, \dots, p_{10} \), so does $p_{12}$, $p_{13}$, and all the way until $p_{n+10}$, giving $n$ people against 10 others, or $10n$ games. In each of these games, the player has a $\frac{1}{2}$ expected chance of winning, giving us $5n$. Combining everything together gives us the above equation.

 And thus we have the system of equations
 \[P = \left( \frac{1}{2} \binom{n+10}{2} \right)(n+10)\]
 \[\frac{1}{2} P = 10 \left( \frac{1}{2} \binom{10}{2} \right) + n \left( \frac{1}{2} \binom{n}{2} \right) + 5n\]

 Through some simple algebra and combinatorial expansions, we find that the values for \( n \) satisfy the equation
 \[(n-15)(n-(8+2\sqrt{31}))(n-(8-2\sqrt{31})) = 0\]

 The three values of \( n \) are
 \[n=15, \quad n=8+2\sqrt{31}, \quad n=8-2\sqrt{31}\]

 However, we require the number of people, and thus the two latter values are obviously incorrect.

 Thus, \( n=15 \). We had \( n+10 \) people in the beginning, and thus the number of people at the tournament were \( 15+10 = \boxed{25} \).

 




	\item (\textit{AIME I 1985}) Three $12$ cm $\times 12$ cm \href{https://artofproblemsolving.com/wiki/index.php?title=Square\_(geometry)}{ squares} are each cut into two pieces $A$ and $B$, as shown in the first figure below, by joining the \href{https://artofproblemsolving.com/wiki/index.php?title=Midpoint}{midpoints} of two adjacent sides. These six pieces are then attached to a \href{https://artofproblemsolving.com/wiki/index.php?title=Regular\_polygon}{ regular} \href{https://artofproblemsolving.com/wiki/index.php?title=Hexagon}{hexagon}, as shown in the second figure, so as to fold into a \href{https://artofproblemsolving.com/wiki/index.php?title=Polyhedron}{polyhedron}. What is the \href{https://artofproblemsolving.com/wiki/index.php?title=Volume}{volume} (in $\mathrm{cm}^3$) of this polyhedron?


 
                    \begin{figure}[H]
                        \centering
                        \includegraphics[width=0.5\textwidth]{Problem/AIME_1985_I_15_Problem_Figure_1.png}
                    \end{figure}
                    


 
 \textbf{Answer: }864



\textbf{Solution 1 (For the visualizers)}

Note that gluing two of the given polyhedra together along a hexagonal face (rotated $60^\circ$ from each other) yields a \href{https://artofproblemsolving.com/wiki/index.php?title=Cube\_(geometry)}{ cube}, so the volume is $\frac12 \cdot 12^3 = 864$, so our answer is $\boxed{864}$.

 Image: 
                    \begin{figure}[H]
                        \centering
                        \includegraphics[width=0.5\textwidth]{Solution/AIME_1985_I_15_Solution_Figure_1.png}
                    \end{figure}
                    

 





\textbf{Solution 2}

Taking a reference to the above left diagram, obviously SP, SQ and SR are perpendicular to each other. Thus expanding plane SPQ, SQR and SRP to intersect with the plane XYZ that contains the regular hexagon, we form a pyramid with S the top vertex and the base being an equilateral triangle with side length of 18 $\sqrt{2}$ . This pyramid has a volume of 972, because it is also 1/6 of the volume of a cube with side length of 18. Then subtracting 3 congruent pyramids with volume of 36 each, we get $864$.

 





\textbf{Solution 3}

Position the polyhedra with the facing hexagon down. From this perspective it is clear that the polyhedra can be thought of as a tetrahedron with 3 smaller tetrahedra removed from each of the three corners of its base. The strategy is then to calculate the volume of the larger tetrahedron, and then subtract the three smaller ones.

 Focusing on the larger tetrahedron, the base is an equilateral triangle formed by extending the sides of the base hexagon; this triangle has sides of length $3\cdot6\sqrt{2}=18\sqrt{2}$. Thus its area is $\frac{\sqrt{3}}{4} \cdot (18 \sqrt{2})^2 = 162\sqrt{3}$.The height of the tetrahedron can be calculated by considering the following right triangle: Leg 1 is the height of the tetrahedron (which we seek); Leg 2 is the distance from a midpoint of one of the sides of the base to the base's center; and the Hypotenuse is the distance from the midpoint of one of the sides of the base to the top vertex. Leg 2 has a length of $\frac{18\sqrt{2}}{2}\cdot\frac{1}{\sqrt{3}} = 3\sqrt{6}$. And the hypotenuse's length is $9\sqrt{2}$ (each non-base face of the tetrahedron is an isosceles right triangle with leg of length 18, whose altitude to its hypotenuse has length $9\sqrt{2}$). Thus, by the Pythagorean theorem the height of the tetrahedron is $6\sqrt{3}$. Altogether, the Volume of the large tetrahedron is $\frac{1}{3}\cdot162\sqrt{3}\cdot6\sqrt{3} = 972$.

 Lastly, each smaller tetrahedron removed from a corner of the base of the larger one has sides whose lengths are $\frac{1}{3}$ that of the larger tetrahedron. Since volume scales with the cube of length, each of these tetrahedra are $\frac{1}{27}$ the volume of the larger one. There are three of them so the final volume is $972 - \frac{3}{27}\cdot972 = 864$.

 




\end{enumerate}